\section{Exact scale transfer and prime dilation of the transport term}
\label{sec:scale-transfer}
%======================================================================

The decomposition $S_n=A_n-T_n$ is exact but static.  This section records the exact \emph{dynamic} content: the transport term $T_n$ is not an independent object but a prime-weighted image of lower-scale Mertens values, transported through the same $2ab$ geometry that governs entry and smoothness.  Every identity here is finite and unconditional; none of them, by itself, is a bound.

\subsection{Source-resolved entry, smoothness, and transport}

For a squarefree source $m=cq$ with $q=\Pplus(m)$, define the entry and transition indices
\begin{equation}
 e(m):=\lfloor\sqrt m\rfloor,
 \qquad
 h(m):=q-1,
 \label{eq:entry-transition-indices}
\end{equation}
so that at stage $n$ the source has entered the prefix when $e(m)\le n$ and is square-root smooth when $h(m)\le n$.  With source weight $\mu(m)$,
\begin{equation}
 \underbrace{\mu(m)\mathbf 1[e(m)\le n]}_{\text{entered}}
 =\underbrace{\mu(m)\mathbf 1[h(m)\le n]}_{\text{smooth}}
 -\underbrace{\bigl(-\mu(m)\bigr)\mathbf 1[e(m)\le n<h(m)]}_{\text{transport}} .
 \label{eq:entered-smooth-transport}
\end{equation}
Summing \cref{eq:entered-smooth-transport} over every source and every finite stage returns the dynamic recombination
\begin{equation}
 S_n=A_n-T_n,
 \label{eq:dynamic-recombination}
\end{equation}
now realized \emph{atom by atom} rather than only in aggregate.

\subsection{Square-root inversion}

Fix the cutoff $R=R_n$ and set
\begin{equation}
 \iota_R(x):=\frac{R^2}{x}.
 \label{eq:sqrt-inversion}
\end{equation}
It is an involution with the single fixed point $R$, and at radius $R=\sqrt{cq}$ it exchanges the two canonical factors, $\iota_R(c)=q$.  It carries the lower interval $[c,R]$ onto the upper interval $[R,R^2/c]$, whose Euclidean lengths obey the exact scaling
\begin{equation}
 \frac{R^2}{c}-R=\frac Rc\,(R-c).
 \label{eq:inversion-length}
\end{equation}
The normalized ordinate is the exact record of this dilation: writing $\lambda_m=\sqrt{q_m/c_m}$,
\begin{equation}
 \frac{q_m^2-c_m^2}{2c_mq_m}
 =\frac{\lambda_m^2-\lambda_m^{-2}}2,
 \qquad
 \log\lambda_m=\tfrac12\operatorname{asinh}\!\left(\frac{q_m^2-c_m^2}{2c_mq_m}\right).
 \label{eq:twoab-dilation}
\end{equation}
Inversion is a geometric map only; it does not preserve primality, and the prime density must re-enter through an explicit multiplier or discrepancy.

\subsection{The exact $2ab$ scale-transfer identity}

Because every sum defining $T_n$ in \cref{eq:Tn} is finite, the cofactor and prime summations may be interchanged.  Using $q>R_n$ and $X_n<R_n^2$, one has $\lfloor X_n/q\rfloor<R_n$, so each prime-first inner sum is a \emph{complete lower-scale Mertens value}.

\begin{theorem}[Prime-first transport transfer]
\label{thm:transport-primefirst}
For every $n\ge1$,
\begin{equation}
 \boxed{
 T_n
 =\sum_{\substack{R_n<q\le X_n\\ q\ \mathrm{prime}}}
 M\!\left(\left\lfloor\frac{X_n}{q}\right\rfloor\right).
 }
 \label{eq:transport-primefirst}
\end{equation}
The complete high-prime transport population is built entirely from Mertens values strictly below the cutoff.
\end{theorem}

Grouping the primes by the reciprocal coordinate $d=\lfloor X_n/q\rfloor$ places the transfer in lower-triangular form.

\begin{theorem}[Reciprocal prime-dilation formula]
\label{thm:prime-dilation}
For every $n\ge1$,
\begin{equation}
 \boxed{
 T_n=\sum_{1\le d<R_n}K_n(d)\,M(d),
 \qquad
 K_n(d)=\pi\!\left(\left\lfloor\frac{X_n}{d}\right\rfloor\right)
 -\pi\!\left(\max\!\left(R_n,\left\lfloor\frac{X_n}{d+1}\right\rfloor\right)\right).
 }
 \label{eq:prime-dilation}
\end{equation}
The kernel $K_n(d)$ is exactly the prime mass in the square-root-inverted interval $X_n/(d+1)<q\le X_n/d$.  The transport term is thus a fixed lower-triangular prime-dilation operator applied to the Mertens sequence.
\end{theorem}

The same construction mirrors the positive-orientation part of the smooth term.  Splitting $A_n=A_n^++A_n^-$ into the canonical pairs $c<q\le R_n$ and the born-smooth remainder $q\le c$, the mirrored part is again an exact prime transform of lower Mertens values,
\begin{equation}
 A_n^+=-\sum_{q\le R_n}M(q-1),
 \label{eq:A-plus-primeform}
\end{equation}
so the difficult cancellation in $S_n=A_n-T_n$ is not between unrelated objects but between two operators generated by the \emph{same} canonical factor geometry.

\subsection{The baseline discrepancy}

For any deterministic baseline $F$ with $F(R)\ne F(c)$, define the scale multiplier and signed discrepancy of an observable $G$ by
\begin{equation}
 \kappa_F(R,c)=\frac{F(R^2/c)-F(R)}{F(R)-F(c)},
 \qquad
 E_{G,F}(R,c)=\bigl[G(R^2/c)-G(R)\bigr]-\kappa_F(R,c)\bigl[G(R)-G(c)\bigr].
 \label{eq:baseline-discrepancy}
\end{equation}
Then identically
\begin{equation}
 G(R^2/c)-G(R)=\kappa_F(R,c)\bigl[G(R)-G(c)\bigr]+E_{G,F}(R,c).
 \label{eq:baseline-split}
\end{equation}
Taking $G=\pi$, weighting by $\mu(c)$, and summing writes transport as a scaled low-prime main term plus one explicit signed discrepancy.  This is the exact location at which a prime-density input must enter: the raw Euclidean multiplier $\kappa_{\mathrm{id}}$ fails at aggregate scale, while the prime-number-theorem, logarithmic-integral, and truncated Riemann-$R$ multipliers remove nearly all of the transport scale, leaving a residual of the same apparent order as $S_n$ itself.

\begin{remark}[Scaling is not a bound]
\label{rem:scaling-not-bound}
Identities \cref{eq:transport-primefirst,eq:prime-dilation,eq:baseline-split} are exact and carry no analytic hypothesis.  They reduce the high-prime transport problem to lower-scale Mertens data, but they do not bound the signed output.  What remains is to prove that the triangular prime-dilation operator is sufficiently contractive on the actual M\"obius sequence, or equivalently to bound the weighted discrepancy $\sum_{c<R_n}\mu(c)E_{\pi,F}(R_n,c)$ together with the scaled main term, with their full signed interaction retained.
\end{remark}

%======================================================================
